\section{Pick 定理：格点、剖分与孔洞}\label{knowledge-pick}
\index{Pick theorem}\index{lattice polygon}\index{Euler characteristic}
简单、非退化、顶点为整数点且无孔的多边形，令 $A$ 为普通非负面积，
$I$ 为严格内部格点数，$B$ 为边界上的不同格点数，则
\[
A=I+\frac B2-1,\qquad I=\frac{|S|-B+2}{2},\qquad
S=\sum_i(x_i y_{i+1}-y_i x_{i+1}).
\]
$S$ 是有向二倍面积，先取绝对值再代入；不要对每一条叉积分别取绝对值。
有向面积接口见第~\ref{compact-polygon_area2}~节（第~\pageref{compact-polygon_area2}~页），
组合用法见第~\ref{usage-example-238}~节（第~\pageref{usage-example-238}~页）。

\paragraph{边界为什么求 gcd}
一条非零整数线段的差向量为 $(u,v)$，令 $g=\gcd(|u|,|v|)$。

所有线段格点恰为起点加 $t(u/g,v/g)$，$t=0,1,\ldots,g$。
这是因为 $(u/g,v/g)$ 互素，任何整数点沿这个方向的参数只能是整数。
每条边只计起点、不计终点，就有 $g$ 个格点；沿简单边界相加，每个顶点恰计一次，故
\[
B=\sum_i\gcd(|x_{i+1}-x_i|,|y_{i+1}-y_i|).
\]
共线的额外边界顶点不改变这个总和。重复走边或自交时不能沿用“每点恰计一次”的论证。

\paragraph{空格点三角形的面积}
这里的“空”指除三个顶点外，内部和边上都没有格点。
平移一个顶点到原点，取一条边 $(u,v)$；它必须是本原向量，即 $\gcd(|u|,|v|)=1$。
由裴蜀取 $ru+sv=1$，整数矩阵
\[
M=\begin{pmatrix}r&s\\-v&u\end{pmatrix},\qquad \det M=1
\]
将该边变成 $(1,0)$，同时双向保留整数格点和面积。
必要时把纵坐标取反，再作整数剪切 $x\gets x-qy$，第三点可写成 $(a,b)$，$0\le a<b$、$b>0$。
若 $b>1$，高度 $y=1$ 的闭截面为 $[a/b,\,1+(a-1)/b]$：
$a=0$ 时含格点 $(0,1)$，$a\ge1$ 时含 $(1,1)$，都是非顶点格点，矛盾。
所以 $b=1$，面积恰为 $1/2$。
仅三条边本原还不够：$(0,0),(2,1),(1,2)$ 的边都本原，但内部有 $(1,1)$。
仅内部无格点也不够：$(0,0),(k,0),(0,1)$ 的面积为 $k/2$，边上可能有其他格点。

\paragraph{相容剖分与 Euler 计数}
先只使用多边形的整数顶点作三角剖分，不引入非整数辅助点，再不断用内部或边上的非顶点格点细分。
点在公共边上时，必须同时细分两侧三角形，不能留下悬挂节点。
每次细分都增加正面积三角形数，而每个三角形的二倍面积是正整数，数量被总二倍面积界住，因此过程终止。
最终每个小三角形面积 $1/2$，所有格点都是剖分顶点，每条边界小边是本原线段。
令 $V,E,F$ 分别为顶点、边、区域内小三角形数（不计无界面），则
\[
V=I+B,\qquad F=2A,\qquad 3F=2E-B,\qquad V-E+F=1.
\]
消去 $V,E,F$ 即得 Pick 公式。这是定理证明；比赛计算不需要实际构造剖分。

\newpage
\paragraph{孔洞、方向与不能套用的情况}
若区域连通，外边界与 $h$ 条孔洞边界都是简单整数多边形，被挖去的孔洞区域严格在外边界内且彼此不交、不接触、不嵌套，
同样的相容剖分满足 $V-E+F=1-h$，所以
\[
A=I+\frac B2-1+h,\qquad I=\frac{2A-B+2}{2}-h.
\]
$I$ 不含孔洞内部及任意边界；$B$ 包含外环与所有内环的边界点。
面积为外环面积减各孔洞面积，或约定外环逆时针、内环顺时针后累加有向面积。
例如外方形 $[0,4]^2$ 去掉孔洞 $(1,3)^2$，有 $A=12,B=24,I=0$；漏掉 $h=1$ 会错得 $I=1$。
这个公式不是本单环用法额外实现了孔洞 API 的声明。

非整数顶点、自交、边界相接或面积退化均不在上述合同内。
蝴蝶结 $(0,0),(2,2),(0,2),(2,0)$ 的有向面积为0，却不能按 Pick 计数。
检查结果非负或分子为偶数，也不能代替简单性检查；正式题是直接保证简单多边形的。

\paragraph{整数位宽与组合用法}
现有整数几何合同为 $n\le10^5$、$|x|,|y|\le C=10^{12}$。
差值和其绝对值可用 long long，但逐项叉积、总二倍面积以及最终分子使用 signed128。
粗界 $|S|\le2nC^2$、$B\le2nC$，在该范围内都安全；内部点数可超过 long long，示例直接打印128位整数。
CSES2193 的正式限制更窄：$3\le n\le10^5$、坐标绝对值不超过 $10^9$，只输入一个简单边界环。
复用 \texttt{polygon\_area2} 与标准库 gcd 即可，不新增一份鞋带公式算法。

参考：\url{https://oi-wiki.org/geometry/pick/}；\url{https://cses.fi/problemset/task/2193}。
